\subsection{Was ist das hier überhaupt?}
\label{sec:was-ist-das}

Dieses Projekt misst, wie sich verschiedene \textbf{JVM-Konfigurationen}
auf den Ressourcenverbrauch und die Performance einer
Spring-Boot-Anwendung in Docker-Containern auswirken. Es ist kein
Produktionsprojekt -- es ist das \textbf{Mess-Werkzeug} für eine
wissenschaftliche Arbeit zum Thema „Ressourcenoptimierung von
javabasierten Containeranwendungen im Cloudbetrieb".

Das Framework besteht aus zwei Teilen:

\begin{enumerate}
	\item
	      \textbf{Die Anwendung unter Test} -- eine Spring-Boot-REST-API, die
	      synthetische Workloads bereitstellt (JSON-Serialisierung,
	      Speicherallokation) sowie einen realen Workload (EBICS-Banktransfer
	      über einen echten Bankserver).
	\item
	      \textbf{Der Benchmark-Runner} -- ein Java-CLI-Programm, das
	      automatisch Docker-Container startet, JVM-Flags injiziert, Metriken
	      erfasst und die Ergebnisse als CSV, JSON und Excel exportiert.
\end{enumerate}

\subsection{Welche JVM-Konfigurationen werden verglichen?}
\label{sec:jvm-konfigurationen}

\subsubsection{Ebene 1: 20 Flag-Analyse-Konfigurationen}
\label{sec:ebene1}

Alle verwenden dasselbe Docker-Image (\texttt{eclipse-temurin:25-jre}).
Unterschiede entstehen nur durch die JVM-Flags.

\begin{longtable}[]{@{}
	>{\raggedright\arraybackslash}p{0.28\linewidth}
	>{\raggedright\arraybackslash}p{0.67\linewidth}@{}}
	\thead{Gruppe}     & \thead{Konfigurationen}                                 \\
	\endhead
	\rowcolor{rowalt}
	Garbage Collectors & G1GC (Default), ZGC, Shenandoah, Parallel GC, Serial GC \\
	G1-Tuning          & MaxGCPauseMillis=50, Heap 256m, Heap 512m               \\
	\rowcolor{rowalt}
	JVM-Interna        & Compressed OOPs aus, Compact Object Headers (JEP 450)   \\
	Cloud-Relevanz     & MaxRAMPercentage=75, TieredStopAtLevel=1 (kein C2)      \\
	\rowcolor{rowalt}
	Kombinationen      & 8 Konfigurationen -- siehe Tabelle unten                \\
\end{longtable}

\noindent Die 8 Flag-Kombinationen im Detail:

\begin{longtable}[]{@{}
	>{\raggedright\arraybackslash}p{0.29\linewidth}
	>{\raggedright\arraybackslash}p{0.33\linewidth}
	>{\raggedright\arraybackslash}p{0.30\linewidth}@{}}
	\thead{Config-Name}           & \thead{JVM-Flags}                                                                                                   & \thead{Hypothese / Ziel}                                                       \\
	\endhead
	\rowcolor{rowalt}
	\texttt{serial-gc-256m}       & \texttt{-XX:+UseSerialGC -Xmx256m}                                                                                  & Minimaler Footprint: Single-Thread-GC + kleiner Heap                           \\
	\texttt{zgc-heap-512m}        & \texttt{-XX:+UseZGC -Xmx512m}                                                                                       & ZGC mit definiertem Heap-Limit, mehr Spielraum für concurrent GC               \\
	\rowcolor{rowalt}
	\texttt{shenandoah-heap-512m} & \texttt{-XX:+UseShenandoahGC -Xmx512m}                                                                              & Shenandoah mit Heap-Limit, Vergleichspunkt zu ZGC                              \\
	\texttt{tiered-stop-1-serial} & \texttt{-XX:TieredStopAtLevel=1 -XX:+UseSerialGC}                                                                   & C1-only + Serial GC: schnellster Start, bester GC auf 1 CPU                    \\
	\rowcolor{rowalt}
	\texttt{g1-coh-on}            & \texttt{-XX:+UseG1GC}\newline\texttt{-XX:+UnlockExperimentalVMOptions}\newline\texttt{-XX:+UseCompactObjectHeaders} & G1 + Compact Object Headers (JEP 450): kleinere Objekte, weniger GC-Druck      \\
	\texttt{parallel-gc-256m}     & \texttt{-XX:+UseParallelGC -Xmx256m}                                                                                & Durchsatz-GC mit kleinem Heap unter Memory-Pressure                            \\
	\rowcolor{rowalt}
	\texttt{g1-large-young}       & \texttt{-XX:+UseG1GC -XX:NewRatio=1}                                                                                & G1 mit 50\% Young Gen (statt \textasciitilde{}33\%), weniger Full GCs erwartet \\
	\texttt{zgc-tiered-stop-1}    & \texttt{-XX:+UseZGC -XX:TieredStopAtLevel=1}                                                                        & ZGC + C1-only: niedrige Pausen mit schnellem Start (Serverless)                \\
\end{longtable}

\subsubsection{Ebene 2: 12 Laufzeitprofile}
\label{sec:ebene2}

Unterschiedliche Docker-Images, unterschiedliche JVM-Implementierungen:

\begin{longtable}[]{@{}
	>{\raggedright\arraybackslash}p{0.20\linewidth}
	>{\raggedright\arraybackslash}p{0.40\linewidth}
	>{\raggedright\arraybackslash}p{0.32\linewidth}@{}}
	\thead{Profil}         & \thead{Image / Base Image}                         & \thead{Was ist anders?}                                                 \\
	\endhead
	\rowcolor{rowalt}
	P01 HotSpot Standard   & \texttt{eclipse-temurin:25-jre}                    & Referenzpunkt (G1 + MaxRAMPercentage=75)                                \\
	P02 Fast Startup       & \texttt{eclipse-temurin:25-jre}                    & TieredStopAtLevel=1 (kein C2)                                           \\
	\rowcolor{rowalt}
	P03 Low Latency        & \texttt{eclipse-temurin:25-jre}                    & ZGC                                                                     \\
	P04 OpenJ9             & \texttt{ibm-semeru-runtimes:open-25-jre}           & Andere JVM-Implementierung (IBM)                                        \\
	\rowcolor{rowalt}
	P05 Native Image       & \texttt{ghcr.io/graalvm/native-image-community:25} & Kein JVM, AOT-kompiliert                                                \\
	P06 OpenJ9 balanced    & \texttt{ibm-semeru-runtimes:open-25-jre}           & \texttt{-Xgcpolicy:balanced} (region-basiert, NUMA-aware)               \\
	\rowcolor{rowalt}
	P07 OpenJ9 optthruput  & \texttt{ibm-semeru-runtimes:open-25-jre}           & \texttt{-Xgcpolicy:optthruput} (paralleler STW-GC, Durchsatz-optimiert) \\
	P08 OpenJ9 optavgpause & \texttt{ibm-semeru-runtimes:open-25-jre}           & \texttt{-Xgcpolicy:optavgpause} (concurrent MS, Pausen-optimiert)       \\
	\rowcolor{rowalt}
	P09 HotSpot 256m       & \texttt{eclipse-temurin:25-jre}                    & Heap begrenzt auf 256 MB                                                \\
	P10 OpenJ9 256m        & \texttt{ibm-semeru-runtimes:open-25-jre}           & Heap 256 MB, direkter JVM-Vergleich                                     \\
	\rowcolor{rowalt}
	P11 CDS                & \texttt{eclipse-temurin:25-jre}                    & Dynamic Class Data Sharing (schnellerer Start)                          \\
	P12 GraalVM JIT        & \texttt{ghcr.io/graalvm/jdk-community:25}          & GraalVM JVMCI statt HotSpot C2                                          \\
\end{longtable}

\subsection{Was wird gemessen?}
\label{sec:metriken}

\subsubsection{Aktuelle Container-Konfiguration (fix für alle Runs)}
\label{sec:container-config}

\begin{longtable}[]{@{}
	>{\raggedright\arraybackslash}p{0.22\linewidth}
	>{\raggedright\arraybackslash}p{0.73\linewidth}@{}}
	\thead{Parameter} & \thead{Wert}                                           \\
	\endhead
	\rowcolor{rowalt}
	CPU-Limit         & \texttt{--cpus 1}                                      \\
	RAM-Limit         & \texttt{--memory 768m}                                 \\
	\rowcolor{rowalt}
	Swap              & deaktiviert (\texttt{--memory-swap 768m})              \\
	Netzwerk          & Host -> Container (lokal, kein echter Netzwerk-Jitter) \\
\end{longtable}

\subsubsection{Gemessene Metriken pro Run}
\label{sec:csv-spalten}

\begin{longtable}[]{@{}
	>{\raggedright\arraybackslash}p{0.18\linewidth}
	>{\raggedright\arraybackslash}p{0.77\linewidth}@{}}
	\thead{Kategorie}  & \thead{Metriken}                                                                                                                               \\
	\endhead
	\rowcolor{rowalt}
	\textbf{Startup}   & \texttt{readinessMs} (Container-Start bis HTTP 200), \texttt{firstSeconds} (erster Workload-Request)                                           \\
	\textbf{Latenz}    & \texttt{latencyP50}, \texttt{latencyP95}, \texttt{latencyP99}, \texttt{latencyMean} (aus 500 Mess-Requests)                                    \\
	\rowcolor{rowalt}
	\textbf{Durchsatz} & \texttt{throughputReqPerSec}, \texttt{totalMeasureTimeSeconds}                                                                                 \\
	\textbf{CPU}       & \texttt{cpuLoadAvg} (Durchschnitt aus 10 docker-stats-Snapshots während Last)                                                                  \\
	\rowcolor{rowalt}
	\textbf{Memory}    & \texttt{memLoadAvg}, \texttt{memLoadMax} (\% des Container-Limits)                                                                             \\
	\textbf{GC}        & \texttt{gcCount}, \texttt{gcFullCount}, \texttt{gcTotalPauseMs}, \texttt{gcMaxPauseMs}, \texttt{gcOverheadPercent}, \texttt{gcPeakHeapAfterMb} \\
	\rowcolor{rowalt}
	\textbf{Metadaten} & \texttt{configName}, \texttt{dockerImage}, \texttt{effectiveJavaToolOptions}, \texttt{runtimeModel}, \texttt{repetition}                       \\
\end{longtable}

Memory-Werte kommen direkt aus \texttt{docker stats} (cgroup-basiert,
also das, was der Kernel dem Container zuteilt -- inkl. JVM Heap,
Metaspace, Thread-Stacks, Code-Cache).

\subsection{Welche Workloads werden simuliert?}
\label{sec:workloads}

\begin{longtable}[]{@{}
	>{\raggedright\arraybackslash}p{0.10\linewidth}
	>{\raggedright\arraybackslash}p{0.25\linewidth}
	>{\raggedright\arraybackslash}p{0.38\linewidth}
	>{\raggedright\arraybackslash}p{0.18\linewidth}@{}}
	\thead{Szenario}      & \thead{Endpoint}                                    & \thead{Was passiert}                                                   & \thead{Workload-Typ}                              \\
	\endhead
	\rowcolor{rowalt}
	\textbf{JSON}         & \texttt{GET /json?n=200000}                         & Erzeugt 200.000 Java-Objekte, serialisiert zu JSON                     & CPU-intensiv, GC-Druck durch kurzlebige Objekte   \\
	\textbf{Alloc}        & {\small\ttfamily GET /alloc?\allowbreak n=10000000} & Allokiert 10 Mio. kurzlebige Byte-Arrays                               & GC-intensiv, stresst den Garbage Collector direkt \\
	\rowcolor{rowalt}
	\textbf{EBICS Upload} & \texttt{GET /ebics/upload}                          & Echte EBICS-Banktransaktion (SEPA-Überweisung) über lokalen Bankserver & I/O + TLS + Krypto, realer Workload               \\
\end{longtable}

Der EBICS-Workload läuft gegen einen echten TravicLink-Bankserver
(Docker Compose: PostgreSQL 16 + Applikationsserver), der lokal läuft.

\subsection{Welche Base-Images sind im Einsatz?}
\label{sec:base-images}

\begin{longtable}[]{@{}
	>{\raggedright\arraybackslash}p{0.40\linewidth}
	>{\raggedright\arraybackslash}p{0.20\linewidth}
	>{\raggedright\arraybackslash}p{0.34\linewidth}@{}}
	\thead{Base Image}                                 & \thead{Typ}               & \thead{Verwendet für}             \\
	\endhead
	\rowcolor{rowalt}
	\texttt{eclipse-temurin:25-jre}                    & HotSpot JRE (Adoptium)    & Standard-JVM-Konfigurationen, CDS \\
	\texttt{ibm-semeru-runtimes:open-25-jre}           & OpenJ9 JRE (IBM/Eclipse)  & Alternative JVM-Implementierung   \\
	\rowcolor{rowalt}
	\texttt{ghcr.io/graalvm/native-image-community:25} & GraalVM (nur Build-Stage) & Native Image Build (Multi-Stage)  \\
	\texttt{ghcr.io/graalvm/jdk-community:25}          & GraalVM JDK               & GraalVM JIT (JVMCI)               \\
	\rowcolor{rowalt}
	\texttt{debian:bookworm-slim}                      & Runtime-Basis             & Native Image Runtime-Stage        \\
\end{longtable}
